% Source-aligned English for printed p.194.
\clearpage
\SourcePage{194}

\LayerTarget{tohoku:setup:g-module-sheaf-notation}
$H^p(G,A)$. Thus $\Gamma^G$ is the functor $h_k$ in the preceding notation,
and

\LayerTarget{tohoku:definition:g-invariant-sections-sheaf}
$\Gamma^G(A)=\Gamma(A)^G$, the group of elements of $\Gamma(A)$ invariant
under $G$, and this equals $\Gamma(A^G)$, where $A^G$ denotes the subsheaf of
$A$ formed by germs of sections invariant under $G$. (If $G$ admits a finite
system of generators, then $A^G$ is also the set of elements of $A$ invariant
under $G$.)

\LayerTarget{tohoku:definition:sheaf-group-cohomology}
The $\mathbf H^p(G,A)$ ($-\infty<p<+\infty$) are sheaf functors of $A$ and
form the derived cohomological functor of
$\boldsymbol{\Gamma}^G(A)=A^G$. There are canonical homomorphisms

\LayerTarget{tohoku:map:group-cohomology-stalk}
\[
\mathbf H^p(G,A)(x)\longrightarrow H^p(G,A(x))
\]

\LayerTarget{tohoku:remark:locality-group-cohomology}
which are bijective in every case that we shall have to consider, and always
when $G$ is finite, by Theorem 4.2.2 or Corollary 1 to Proposition 4.2.4, as
one prefers! These functors are local in nature; that is, they commute with
restriction to an open set.

\LayerTarget{tohoku:definition:equivariant-cohomology}
The $H^p(X;G,A)$ form the derived cohomological functor of $\Gamma^G(A)$.
We have the following theorem.

\LayerTarget{tohoku:theorem:4.4.1}
\noindent\textsc{Theorem 4.4.1.} \textit{There exist two cohomological spectral
functors on the category of $G$-sheaves, abutting to the graded functor
$(H^m(X;G,A))$, whose initial terms are respectively}

\LayerTarget{tohoku:eq:4.4.1}
\[
\tag{4.4.1}
\begin{aligned}
\mathrm I_2^{p,q}(A)&=H^p(X,\mathbf H^q(G,A)),\\
\mathrm {II}_2^{p,q}&=H^p(G,H^q(X,A)).
\end{aligned}
\]

\LayerTarget{tohoku:construction:operator-complex-spectral-sequences}
\textit{These are also the two spectral sequences of the ``complex with
operators'' $C(A)$, where $C(A)$ is the canonical resolution of $A$, or the
two hypercohomology spectral sequences of $X$ with respect to the complex
$C(G,A)$ of ``$G$-cochains with coefficients in the sheaf $A$''. These two
spectral sequences are trivial if $G$ acts trivially on $A$ and, in addition,
$A$ is a vector space over the field $k$; then there is a canonical
isomorphism}

\LayerTarget{tohoku:eq:4.4.2}
\[
\tag{4.4.2}
H^n(X;G,A)=
\bigoplus_{p+q=n}
\mathrm{Hom}_k(H_p(G,k),H^q(X,A)).
\]

\textit{More generally, suppose only that $G$ acts trivially on $A$. Then
Corollary 2 to Proposition 4.3.2 gives a canonical exact sequence}

\LayerTarget{tohoku:eq:4.4.3}
\[
\tag{4.4.3}
\begin{aligned}
0\longrightarrow{}&
\bigoplus_{p+q=n-1}
\mathrm{Ext}_{\mathbf Z}^1(H_p(G,\mathbf Z),H^q(X,A))\\
\longrightarrow{}&H^n(X;G,A)
\longrightarrow
\bigoplus_{p+q=n}
\mathrm{Hom}_{\mathbf Z}(H_p(G,\mathbf Z),H^q(X,A))
\longrightarrow0.
\end{aligned}
\]

\LayerTarget{tohoku:remark:4.4.1:1}
\noindent\textsc{Remarks.} 1. In this last formula we had implicitly taken
$k=\mathbf Z$. This is immaterial, however, because if $k$ is arbitrary and
$A$ is a sheaf of $k$-modules on which $G$ acts, then $H^p(G,A)$ and
$H^p(X;G,A)$ are the same whether $A$ is regarded as a sheaf of
$k(G)$-modules or as a sheaf of $\mathbf Z(G)$-modules.

\LayerTarget{tohoku:proof:remark-4.4.1:1}
To see this, it suffices to prove that if $A$ is an injective sheaf of
$k(G)$-modules, then $\mathrm{Ext}_{\mathbf Z(G)}^n(X;\mathbf Z,A)$ and
$\mathbf{Ext}_{\mathbf Z(G)}^n(\mathbf Z,A)$, where $\mathbf Z$ and
$\mathbf Z(G)$ denote the constant sheaves that they define, vanish for
$n>0$. For the first, this follows from spectral sequence $II$, which shows
that it is isomorphic to $H^n(G,\Gamma(A))$. This vanishes because $\Gamma(A)$
is an injective $k$-module, for which the result invoked is well known.

\TohokuContinuation{195}
